%% trilogy-style.tex -- typographic design for the submission.
%%
%% Loaded by main.tex with %% trilogy-style.tex -- typographic design for the submission.
%%
%% Loaded by main.tex with %% trilogy-style.tex -- typographic design for the submission.
%%
%% Loaded by main.tex with %% trilogy-style.tex -- typographic design for the submission.
%%
%% Loaded by main.tex with \input{shared/trilogy-style}, after the class, fontenc and
%% inputenc, the AMS packages, natbib or biblatex, and imakeidx, and before the document's
%% theorem declarations and macros. The title, theorem names, local notes, and short
%% running heads are defined in main.tex.
%%
%% Principles: readability first, then understated elegance. The document is read mostly
%% on screen, so the page is set for that: 12pt type on a full-width letter page with 1in
%% side margins, one-sided, with the same quiet head on every page. Black text only; one
%% readable face with matching mathematics; headings distinguished by size, weight and
%% space rather than colour or rules; statements in the standard amsthm manner; classic
%% title page. The main.tex uses \documentclass[12pt,oneside,
%% openany]{book}.

\makeatletter

%% ---------------------------------------------------------------------------
%% Type: Palatino for text and mathematics (mathpazo, with true small capitals).
%% Its larger x-height and sturdier strokes read better than Latin Modern, on screen
%% and in dense formulas, and the math italic and symbols are drawn to match the text.
%% Latin Modern supplies the sans-serif and typewriter faces (URLs). bm comes after
%% the math fonts so that \bm finds bold Palatino.
%% ---------------------------------------------------------------------------
\usepackage{lmodern}
\usepackage[sc]{mathpazo}
\usepackage{bm}
\linespread{1.1}
\usepackage{microtype}
% Font expansion limited to 1.5% (default 2%): steadier colour, and it avoids the
% hairline overfull lines that maximal shrink plus protrusion can leave.
\microtypesetup{stretch=15,shrink=15}

%% ---------------------------------------------------------------------------
%% Page, set for reading on screen: US letter, 1in side margins (a 6.5in line, about
%% 85 characters of Palatino 12pt), and slightly smaller top and bottom margins with
%% the head and foot inside them, so a page fitted to the window shows large type.
%% ---------------------------------------------------------------------------
\usepackage{geometry}
\geometry{letterpaper,left=1in,right=1in,top=0.75in,bottom=0.75in,includeheadfoot,
  headheight=15pt,headsep=16pt,footskip=26pt}
\setlength{\parindent}{1.5em}
\setlength{\parskip}{0pt}
\setlength{\emergencystretch}{3em}
% A paragraph may set a slightly loose line rather than push a word into the margin.
\tolerance=1000
\clubpenalty=3000 \widowpenalty=3000 \displaywidowpenalty=1500
\raggedbottom
\allowdisplaybreaks[2]

%% Lists: compact, indented like a paragraph.
\usepackage{enumitem}
\setlist{topsep=4pt plus 2pt minus 1pt,itemsep=2pt plus 1pt,parsep=0pt}
\setlist[itemize]{leftmargin=1.5em}
\setlist[enumerate]{leftmargin=2em}

\usepackage[newparttoc]{titlesec}
\usepackage{titletoc,fancyhdr}

%% Letterspaced small capitals, for labels such as "Chapter 3".
\newcommand{\trilogy@spacedsc}[1]{{\scshape\textls[90]{\MakeLowercase{#1}}}}
%% Headings are ragged right and unhyphenated; a title that needs two lines is broken
%% into lines of similar length rather than leaving a single word on the second line.
\newcommand{\trilogy@headpar}{\rightskip=0pt plus 0.5\hsize\relax
  \parfillskip=0pt plus 0.5\hsize\relax\hyphenpenalty=10000\relax\exhyphenpenalty=10000\relax}

%% ---------------------------------------------------------------------------
%% Headings (titlesec): a regular-weight chapter title under a spaced small-capitals
%% label; bold section titles; plain part pages; no colour and no rules.
%% ---------------------------------------------------------------------------
\titleformat{\part}[display]
  {\normalfont\centering}
  {\large\trilogy@spacedsc{\partname\ \thepart}}
  {22pt}
  {\fontsize{22}{28}\selectfont}
\assignpagestyle{\part}{empty}

\titleformat{\chapter}[display]
  {\normalfont\trilogy@headpar}
  {\trilogy@spacedsc{\chaptertitlename\ \thechapter}}
  {12pt}
  {\fontsize{20}{25}\selectfont}
\titleformat{name=\chapter,numberless}[display]
  {\normalfont\trilogy@headpar}
  {}
  {0pt}
  {\fontsize{20}{25}\selectfont\trilogy@starmark}
\titlespacing*{\chapter}{0pt}{12pt}{30pt}
% An unnumbered chapter (front matter, contents, references, index) sets both running heads.
\newcommand{\trilogy@starmark}[1]{#1\markboth{#1}{#1}}

\titleformat{\section}
  {\normalfont\fontsize{13.5}{17}\selectfont\bfseries\trilogy@headpar}{\thesection}{0.75em}{}
\titlespacing*{\section}{0pt}{22pt plus 4pt minus 2pt}{9pt plus 2pt}
\titleformat{\subsection}
  {\normalfont\normalsize\bfseries\trilogy@headpar}{\thesubsection}{0.75em}{}
\titlespacing*{\subsection}{0pt}{15pt plus 3pt minus 2pt}{6pt plus 1pt}
\titleformat{\paragraph}[runin]
  {\normalfont\normalsize\bfseries}{}{0pt}{}
\titlespacing*{\paragraph}{0pt}{9pt plus 2pt minus 1pt}{0.6em}

\setcounter{tocdepth}{1}
\setcounter{secnumdepth}{1}

%% ---------------------------------------------------------------------------
%% Running heads (one-sided, the same on every page): the chapter number and title at
%% the left, and at the right the section number and title followed by the page number,
%% all in small italic except the numbers, with no rule. When chapter and section titles
%% do not both fit on the line, the right head shows only the section number (§1.2). Chapter-
%% opening pages carry only a centred page number at the foot; part pages carry nothing.
%%
%% A long title can be given a short form in main.tex:
%%   \shorthead{<title exactly as in the source>}{<short form>}
%% A chapter head too wide for the line falls back to "Chapter N" (with a warning).
%% ---------------------------------------------------------------------------
\newcommand{\shorthead}[2]{\expandafter\def\csname trilogy@sh@\detokenize{#1}\endcsname{#2}}
\newcommand{\trilogy@lookup}[1]{%
  \ifcsname trilogy@sh@\detokenize{#1}\endcsname
    \expandafter\let\expandafter\trilogy@headtext\csname trilogy@sh@\detokenize{#1}\endcsname
  \else
    \def\trilogy@headtext{#1}%
  \fi}
\newcommand{\trilogy@headfont}{\normalfont\small\itshape}
% A mark is \trilogy@hd{number}{title}; the number is empty for unnumbered chapters.
\DeclareRobustCommand{\trilogy@hd}[2]{%
  \ifx\relax#1\relax\else{\upshape#1}\hspace{0.8em}\fi#2}
\newsavebox{\trilogy@chapbox}
\newsavebox{\trilogy@secbox}
\newsavebox{\trilogy@cachedchapbox}
\newsavebox{\trilogy@cachedsecbox}
\newif\iftrilogy@headtoolong
\newif\iftrilogy@cachedheadtoolong
\let\trilogy@cachedheadkey\relax
% Compute the two boxes only when their marks or dimensions change.
\newcommand{\trilogy@computeheads}{%
  \trilogy@headtoolongfalse
  \sbox\trilogy@chapbox{\trilogy@headfont\nouppercase{\leftmark}}%
  \sbox\trilogy@secbox{}%
  \protected@edef\trilogy@l{\leftmark}%
  \protected@edef\trilogy@r{\rightmark}%
  \ifx\trilogy@r\@empty\else\ifx\trilogy@r\trilogy@l\else
    \sbox\trilogy@secbox{\trilogy@headfont\nouppercase{\rightmark}}%
  \fi\fi
  \dimen@=\wd\trilogy@chapbox\advance\dimen@ 4em\relax
  \ifdim\dimen@>\headwidth
    \sbox\trilogy@chapbox{\trilogy@headfont\@chapapp\ \thechapter}%
    \trilogy@headtoolongtrue
  \fi
  \dimen@=\wd\trilogy@chapbox\advance\dimen@\wd\trilogy@secbox\advance\dimen@ 8em\relax
  \ifdim\dimen@>\headwidth
    \sbox\trilogy@secbox{\trilogy@headfont
      \expandafter\let\csname trilogy@hd \endcsname\trilogy@hdnum\nouppercase{\rightmark}}%
  \fi}
% Both fancyhdr fields, and successive pages with the same marks, reuse these
% boxes. The cache is global because fancyhdr evaluates its fields in groups;
% the page number and page-specific warnings remain outside the cached boxes.
\newcommand{\trilogy@heads}[1]{%
  \protected@edef\trilogy@headkey{%
    {\leftmark}{\rightmark}{\the\headwidth}%
    {\fontname\font}{\the\dimexpr1em\relax}{\thechapter}{\@chapapp}}%
  \ifx\trilogy@headkey\trilogy@cachedheadkey\else
    \trilogy@computeheads
    \global\setbox\trilogy@cachedchapbox=\copy\trilogy@chapbox
    \global\setbox\trilogy@cachedsecbox=\copy\trilogy@secbox
    \global\let\trilogy@cachedheadkey\trilogy@headkey
    \iftrilogy@headtoolong
      \global\trilogy@cachedheadtoolongtrue
    \else
      \global\trilogy@cachedheadtoolongfalse
    \fi
  \fi
  \setbox\trilogy@chapbox=\copy\trilogy@cachedchapbox
  \setbox\trilogy@secbox=\copy\trilogy@cachedsecbox
  \iftrilogy@cachedheadtoolong
    \ifx\relax#1\relax\else\PackageWarningNoLine{trilogy-style}{Running head too long on
      page \thepage; give the chapter a \string\shorthead}\fi
  \fi}
% The section number alone, for a right head with no room for the section title. (A mark
% holds the robust command's inner macro "\trilogy@hd ", which is what is replaced above.)
\newcommand{\trilogy@hdnum}[2]{\ifx\relax#1\relax\else\S\,{\upshape#1}\fi}

\pagestyle{fancy}
\fancyhf{}
\fancyhead[L]{\trilogy@heads{warn}\usebox\trilogy@chapbox}
\fancyhead[R]{\trilogy@heads{}%
  \ifdim\wd\trilogy@secbox>\z@\usebox\trilogy@secbox\hspace{2em}\fi
  \normalfont\small\thepage}
\renewcommand{\headrulewidth}{0pt}
\renewcommand{\footrulewidth}{0pt}
\fancypagestyle{plain}{\fancyhf{}\fancyfoot[C]{\normalfont\small\thepage}%
  \renewcommand{\headrulewidth}{0pt}\renewcommand{\footrulewidth}{0pt}}
\renewcommand{\chaptermark}[1]{\trilogy@lookup{#1}%
  \markboth{\trilogy@hd{\ifnum\c@secnumdepth>\m@ne\if@mainmatter\thechapter\fi\fi}{\trilogy@headtext}}{}}
\renewcommand{\sectionmark}[1]{\trilogy@lookup{#1}%
  \markright{\trilogy@hd{\ifnum\c@secnumdepth>\z@\thesection\fi}{\trilogy@headtext}}}

%% ---------------------------------------------------------------------------
%% Table of contents (titletoc): parts in spaced small capitals, chapters in roman with
%% hanging numbers, sections indented and smaller with sparse leaders.
%% ---------------------------------------------------------------------------
\contentsmargin{2.2em}
\titlecontents{part}[0pt]
  {\addvspace{16pt plus 2pt}\normalfont}
  {\trilogy@spacedsc{\partname\ \thecontentslabel}\hspace{1em}}
  {}
  {\hfill\contentspage}
  [\addvspace{4pt}]
\titlecontents{chapter}[2.2em]
  {\addvspace{7pt plus 1pt}\normalfont}
  {\contentslabel{2.2em}}
  {\hspace*{-2.2em}}
  {\hfill\contentspage}
\titlecontents{section}[4.8em]
  {\normalfont\small}
  {\contentslabel{2.6em}}
  {\hspace*{-2.6em}}
  {\titlerule*[0.9em]{.}\contentspage}

%% ---------------------------------------------------------------------------
%% Index and bibliography: the index in two columns of small type; the bibliography
%% is "References", listed in the contents.
%% ---------------------------------------------------------------------------
% The last index page is balanced by multicol, which by default accepts a balanced column up
% to 12pt (a line) too long, so that one added entry can push a line into the bottom margin.
% With no overflow allowed, multicol fills that page normally and balances the rest on the
% next page instead; ragged columns keep balanced columns at their natural height.
\indexsetup{othercode={\small\raggedcolumns\maxbalancingoverflow=\z@}}
\@ifpackageloaded{natbib}{%
  \renewcommand{\bibsection}{\chapter*{References}%
    \addcontentsline{toc}{chapter}{References}}%
}{}

%% ---------------------------------------------------------------------------
%% Statements, examples and notes (amsthm). The three standard styles are redefined
%% with one generous spacing, so \theoremstyle{plain|definition|remark} in
%% main.tex picks up the design:
%%   plain       bold label, italic body    (Statement, Theorem, Lemma, ...)
%%   definition  bold label, upright body   (Definition, Example, Technical walkthrough)
%%   remark      italic label, upright body (Caution, Seminar question)
%% ---------------------------------------------------------------------------
\newlength{\trilogyenvskip}
\setlength{\trilogyenvskip}{10pt plus 3pt minus 2pt}
\newtheoremstyle{plain}{\trilogyenvskip}{\trilogyenvskip}
  {\itshape}{}{\bfseries}{.}{0.5em}{}
\newtheoremstyle{definition}{\trilogyenvskip}{\trilogyenvskip}
  {\normalfont}{}{\bfseries}{.}{0.5em}{}
\newtheoremstyle{remark}{\trilogyenvskip}{\trilogyenvskip}
  {\normalfont}{}{\itshape}{.}{0.5em}{}

%% Run-in notes ("Used in:", "Sources.", ...): an italic label and ordinary text,
%% a little space above. \trilogynote{<label>}{<text>}
\newcommand{\trilogynote}[2]{\par\addvspace{5pt plus 1pt}\noindent{\itshape #1}\enspace #2\par}

%% "Used in:" line after a statement: \usedin{<references>}
\newcommand{\usedin}[1]{\trilogynote{Used in:}{#1}}

%% Technical walkthrough: a numbered, multi-paragraph guided calculation,
%% \begin{walkthrough}[<optional title>] ... \end{walkthrough}. Definition style;
%% it ends with an open square at the right margin, like a proof, so the reader sees
%% where it stops. The name may be changed with \renewcommand{\walkthroughname}{...}.
\newcommand{\walkthroughname}{Technical walkthrough}
\newcommand{\walkthroughend}{\ensuremath{\square}}
%% Set-off statement: definition style, bold label and upright text, with the whole statement
%% indented 1.5em on both sides like the asides, instead of a box. Displays stay centred.
%% Declare with \theoremstyle{trilogysetoff}\newtheorem{<env>}{<name>}[chapter].
\newtheoremstyle{trilogysetoff}{\trilogyenvskip}{\trilogyenvskip}
  {\normalfont\leftskip=1.5em\rightskip=1.5em\relax}{}{\bfseries}{.}{0.5em}{}

\theoremstyle{definition}
\newtheorem{trilogywalk}{\walkthroughname}[chapter]
\newenvironment{walkthrough}[1][]
  {\ifx\relax#1\relax\begin{trilogywalk}\else\begin{trilogywalk}[#1]\fi
   \pushQED{\qed}\renewcommand{\qedsymbol}{\walkthroughend}}
  {\popQED\end{trilogywalk}}
\theoremstyle{plain}

%% Set-off asides, \begin{sayit} ... \end{sayit} ("How to say it") and
%% \begin{fiveminute} ... \end{fiveminute} ("Five-minute explanation"): a spaced
%% small-capitals label and a body indented on both sides; no box, no rule, no tint.
%% Inside sayit, \tophysicist and \tomathematician start the two voices.
\newenvironment{trilogyaside}[1]
  {\par\addvspace{11pt plus 3pt minus 2pt}%
   \list{}{\leftmargin=1.5em\rightmargin=1.5em\topsep=0pt\partopsep=0pt\parsep=0pt
     \listparindent=1.5em\itemsep=0pt}%
   \item\relax{\trilogy@spacedsc{#1}}\par\nopagebreak\vspace{3pt}\noindent\ignorespaces}
  {\endlist\par\addvspace{11pt plus 3pt minus 2pt}}
\newenvironment{sayit}{\begin{trilogyaside}{How to say it}}{\end{trilogyaside}}
\newenvironment{fiveminute}{\begin{trilogyaside}{Five-minute explanation}}{\end{trilogyaside}}
\newcommand{\tophysicist}{\par\addvspace{2pt}\noindent{\itshape To a physicist.}\enspace\ignorespaces}
\newcommand{\tomathematician}{\par\addvspace{4pt}\noindent{\itshape To a mathematician.}\enspace\ignorespaces}

%% Keeping a heading with what follows it: \trilogyneedspace{<length>} starts a new page at
%% once unless at least <length> of the current page remains. The test is made on the spot,
%% so it also works before a longtable (a breakpoint left for later would be taken under
%% longtable's output routine, which then repeats the table head above the heading).
\newcommand{\trilogyneedspace}[1]{\par\penalty-100\begingroup
  \setlength{\dimen@}{#1}\dimen@ii\pagegoal\advance\dimen@ii-\pagetotal
  \ifdim\dimen@>\dimen@ii\ifdim\dimen@ii>\z@\vfil\fi\break\fi\endgroup}

%% Captions: small type (11pt at the 12pt base), "Table A.1." with a full stop.
\usepackage[font=small,labelsep=period,skip=8pt]{caption}

%% Footnotes in \small rather than \footnotesize, for legibility on screen.
\usepackage{etoolbox}
\patchcmd{\@footnotetext}{\footnotesize}{\small}{}{}

%% ---------------------------------------------------------------------------
%% Wide reference tables. \widetablesetup, issued inside a group just before a
%% longtable, lets that table run \trilogy@overhang into each margin, centred on the
%% page: within the group \linewidth is the measure plus twice the overhang, so
%% p{<fraction>\linewidth} columns scale with it. With the 6.5in screen measure the
%% overhang is zero; the command is kept so that tables using it need no change.
%% ---------------------------------------------------------------------------
\newcommand{\trilogy@overhang}{0pt}
\newcommand{\widetablesetup}{%
  \setlength{\LTleft}{-\trilogy@overhang plus 1fill}%
  \setlength{\LTright}{-\trilogy@overhang plus 1fill}%
  \setlength{\linewidth}{\dimexpr\textwidth+2\dimexpr\trilogy@overhang\relax\relax}}

%% ---------------------------------------------------------------------------
%% Links: black and unboxed; they still work. No author in the PDF metadata.
%% The PDF opens on the document alone (no side panel), at full page width,
%% with the document title (not the file name) in the window bar.
%% ---------------------------------------------------------------------------
\usepackage[hidelinks,bookmarksnumbered=true]{hyperref}
\usepackage{bookmark}
\hypersetup{pdfauthor={},pdfpagemode=UseNone,pdfdisplaydoctitle=true,
  pdfstartview={FitH \hypercalcbp{\paperheight}}}

\makeatother
, after the class, fontenc and
%% inputenc, the AMS packages, natbib or biblatex, and imakeidx, and before the document's
%% theorem declarations and macros. The title, theorem names, local notes, and short
%% running heads are defined in main.tex.
%%
%% Principles: readability first, then understated elegance. The document is read mostly
%% on screen, so the page is set for that: 12pt type on a full-width letter page with 1in
%% side margins, one-sided, with the same quiet head on every page. Black text only; one
%% readable face with matching mathematics; headings distinguished by size, weight and
%% space rather than colour or rules; statements in the standard amsthm manner; classic
%% title page. The main.tex uses \documentclass[12pt,oneside,
%% openany]{book}.

\makeatletter

%% ---------------------------------------------------------------------------
%% Type: Palatino for text and mathematics (mathpazo, with true small capitals).
%% Its larger x-height and sturdier strokes read better than Latin Modern, on screen
%% and in dense formulas, and the math italic and symbols are drawn to match the text.
%% Latin Modern supplies the sans-serif and typewriter faces (URLs). bm comes after
%% the math fonts so that \bm finds bold Palatino.
%% ---------------------------------------------------------------------------
\usepackage{lmodern}
\usepackage[sc]{mathpazo}
\usepackage{bm}
\linespread{1.1}
\usepackage{microtype}
% Font expansion limited to 1.5% (default 2%): steadier colour, and it avoids the
% hairline overfull lines that maximal shrink plus protrusion can leave.
\microtypesetup{stretch=15,shrink=15}

%% ---------------------------------------------------------------------------
%% Page, set for reading on screen: US letter, 1in side margins (a 6.5in line, about
%% 85 characters of Palatino 12pt), and slightly smaller top and bottom margins with
%% the head and foot inside them, so a page fitted to the window shows large type.
%% ---------------------------------------------------------------------------
\usepackage{geometry}
\geometry{letterpaper,left=1in,right=1in,top=0.75in,bottom=0.75in,includeheadfoot,
  headheight=15pt,headsep=16pt,footskip=26pt}
\setlength{\parindent}{1.5em}
\setlength{\parskip}{0pt}
\setlength{\emergencystretch}{3em}
% A paragraph may set a slightly loose line rather than push a word into the margin.
\tolerance=1000
\clubpenalty=3000 \widowpenalty=3000 \displaywidowpenalty=1500
\raggedbottom
\allowdisplaybreaks[2]

%% Lists: compact, indented like a paragraph.
\usepackage{enumitem}
\setlist{topsep=4pt plus 2pt minus 1pt,itemsep=2pt plus 1pt,parsep=0pt}
\setlist[itemize]{leftmargin=1.5em}
\setlist[enumerate]{leftmargin=2em}

\usepackage[newparttoc]{titlesec}
\usepackage{titletoc,fancyhdr}

%% Letterspaced small capitals, for labels such as "Chapter 3".
\newcommand{\trilogy@spacedsc}[1]{{\scshape\textls[90]{\MakeLowercase{#1}}}}
%% Headings are ragged right and unhyphenated; a title that needs two lines is broken
%% into lines of similar length rather than leaving a single word on the second line.
\newcommand{\trilogy@headpar}{\rightskip=0pt plus 0.5\hsize\relax
  \parfillskip=0pt plus 0.5\hsize\relax\hyphenpenalty=10000\relax\exhyphenpenalty=10000\relax}

%% ---------------------------------------------------------------------------
%% Headings (titlesec): a regular-weight chapter title under a spaced small-capitals
%% label; bold section titles; plain part pages; no colour and no rules.
%% ---------------------------------------------------------------------------
\titleformat{\part}[display]
  {\normalfont\centering}
  {\large\trilogy@spacedsc{\partname\ \thepart}}
  {22pt}
  {\fontsize{22}{28}\selectfont}
\assignpagestyle{\part}{empty}

\titleformat{\chapter}[display]
  {\normalfont\trilogy@headpar}
  {\trilogy@spacedsc{\chaptertitlename\ \thechapter}}
  {12pt}
  {\fontsize{20}{25}\selectfont}
\titleformat{name=\chapter,numberless}[display]
  {\normalfont\trilogy@headpar}
  {}
  {0pt}
  {\fontsize{20}{25}\selectfont\trilogy@starmark}
\titlespacing*{\chapter}{0pt}{12pt}{30pt}
% An unnumbered chapter (front matter, contents, references, index) sets both running heads.
\newcommand{\trilogy@starmark}[1]{#1\markboth{#1}{#1}}

\titleformat{\section}
  {\normalfont\fontsize{13.5}{17}\selectfont\bfseries\trilogy@headpar}{\thesection}{0.75em}{}
\titlespacing*{\section}{0pt}{22pt plus 4pt minus 2pt}{9pt plus 2pt}
\titleformat{\subsection}
  {\normalfont\normalsize\bfseries\trilogy@headpar}{\thesubsection}{0.75em}{}
\titlespacing*{\subsection}{0pt}{15pt plus 3pt minus 2pt}{6pt plus 1pt}
\titleformat{\paragraph}[runin]
  {\normalfont\normalsize\bfseries}{}{0pt}{}
\titlespacing*{\paragraph}{0pt}{9pt plus 2pt minus 1pt}{0.6em}

\setcounter{tocdepth}{1}
\setcounter{secnumdepth}{1}

%% ---------------------------------------------------------------------------
%% Running heads (one-sided, the same on every page): the chapter number and title at
%% the left, and at the right the section number and title followed by the page number,
%% all in small italic except the numbers, with no rule. When chapter and section titles
%% do not both fit on the line, the right head shows only the section number (§1.2). Chapter-
%% opening pages carry only a centred page number at the foot; part pages carry nothing.
%%
%% A long title can be given a short form in main.tex:
%%   \shorthead{<title exactly as in the source>}{<short form>}
%% A chapter head too wide for the line falls back to "Chapter N" (with a warning).
%% ---------------------------------------------------------------------------
\newcommand{\shorthead}[2]{\expandafter\def\csname trilogy@sh@\detokenize{#1}\endcsname{#2}}
\newcommand{\trilogy@lookup}[1]{%
  \ifcsname trilogy@sh@\detokenize{#1}\endcsname
    \expandafter\let\expandafter\trilogy@headtext\csname trilogy@sh@\detokenize{#1}\endcsname
  \else
    \def\trilogy@headtext{#1}%
  \fi}
\newcommand{\trilogy@headfont}{\normalfont\small\itshape}
% A mark is \trilogy@hd{number}{title}; the number is empty for unnumbered chapters.
\DeclareRobustCommand{\trilogy@hd}[2]{%
  \ifx\relax#1\relax\else{\upshape#1}\hspace{0.8em}\fi#2}
\newsavebox{\trilogy@chapbox}
\newsavebox{\trilogy@secbox}
\newsavebox{\trilogy@cachedchapbox}
\newsavebox{\trilogy@cachedsecbox}
\newif\iftrilogy@headtoolong
\newif\iftrilogy@cachedheadtoolong
\let\trilogy@cachedheadkey\relax
% Compute the two boxes only when their marks or dimensions change.
\newcommand{\trilogy@computeheads}{%
  \trilogy@headtoolongfalse
  \sbox\trilogy@chapbox{\trilogy@headfont\nouppercase{\leftmark}}%
  \sbox\trilogy@secbox{}%
  \protected@edef\trilogy@l{\leftmark}%
  \protected@edef\trilogy@r{\rightmark}%
  \ifx\trilogy@r\@empty\else\ifx\trilogy@r\trilogy@l\else
    \sbox\trilogy@secbox{\trilogy@headfont\nouppercase{\rightmark}}%
  \fi\fi
  \dimen@=\wd\trilogy@chapbox\advance\dimen@ 4em\relax
  \ifdim\dimen@>\headwidth
    \sbox\trilogy@chapbox{\trilogy@headfont\@chapapp\ \thechapter}%
    \trilogy@headtoolongtrue
  \fi
  \dimen@=\wd\trilogy@chapbox\advance\dimen@\wd\trilogy@secbox\advance\dimen@ 8em\relax
  \ifdim\dimen@>\headwidth
    \sbox\trilogy@secbox{\trilogy@headfont
      \expandafter\let\csname trilogy@hd \endcsname\trilogy@hdnum\nouppercase{\rightmark}}%
  \fi}
% Both fancyhdr fields, and successive pages with the same marks, reuse these
% boxes. The cache is global because fancyhdr evaluates its fields in groups;
% the page number and page-specific warnings remain outside the cached boxes.
\newcommand{\trilogy@heads}[1]{%
  \protected@edef\trilogy@headkey{%
    {\leftmark}{\rightmark}{\the\headwidth}%
    {\fontname\font}{\the\dimexpr1em\relax}{\thechapter}{\@chapapp}}%
  \ifx\trilogy@headkey\trilogy@cachedheadkey\else
    \trilogy@computeheads
    \global\setbox\trilogy@cachedchapbox=\copy\trilogy@chapbox
    \global\setbox\trilogy@cachedsecbox=\copy\trilogy@secbox
    \global\let\trilogy@cachedheadkey\trilogy@headkey
    \iftrilogy@headtoolong
      \global\trilogy@cachedheadtoolongtrue
    \else
      \global\trilogy@cachedheadtoolongfalse
    \fi
  \fi
  \setbox\trilogy@chapbox=\copy\trilogy@cachedchapbox
  \setbox\trilogy@secbox=\copy\trilogy@cachedsecbox
  \iftrilogy@cachedheadtoolong
    \ifx\relax#1\relax\else\PackageWarningNoLine{trilogy-style}{Running head too long on
      page \thepage; give the chapter a \string\shorthead}\fi
  \fi}
% The section number alone, for a right head with no room for the section title. (A mark
% holds the robust command's inner macro "\trilogy@hd ", which is what is replaced above.)
\newcommand{\trilogy@hdnum}[2]{\ifx\relax#1\relax\else\S\,{\upshape#1}\fi}

\pagestyle{fancy}
\fancyhf{}
\fancyhead[L]{\trilogy@heads{warn}\usebox\trilogy@chapbox}
\fancyhead[R]{\trilogy@heads{}%
  \ifdim\wd\trilogy@secbox>\z@\usebox\trilogy@secbox\hspace{2em}\fi
  \normalfont\small\thepage}
\renewcommand{\headrulewidth}{0pt}
\renewcommand{\footrulewidth}{0pt}
\fancypagestyle{plain}{\fancyhf{}\fancyfoot[C]{\normalfont\small\thepage}%
  \renewcommand{\headrulewidth}{0pt}\renewcommand{\footrulewidth}{0pt}}
\renewcommand{\chaptermark}[1]{\trilogy@lookup{#1}%
  \markboth{\trilogy@hd{\ifnum\c@secnumdepth>\m@ne\if@mainmatter\thechapter\fi\fi}{\trilogy@headtext}}{}}
\renewcommand{\sectionmark}[1]{\trilogy@lookup{#1}%
  \markright{\trilogy@hd{\ifnum\c@secnumdepth>\z@\thesection\fi}{\trilogy@headtext}}}

%% ---------------------------------------------------------------------------
%% Table of contents (titletoc): parts in spaced small capitals, chapters in roman with
%% hanging numbers, sections indented and smaller with sparse leaders.
%% ---------------------------------------------------------------------------
\contentsmargin{2.2em}
\titlecontents{part}[0pt]
  {\addvspace{16pt plus 2pt}\normalfont}
  {\trilogy@spacedsc{\partname\ \thecontentslabel}\hspace{1em}}
  {}
  {\hfill\contentspage}
  [\addvspace{4pt}]
\titlecontents{chapter}[2.2em]
  {\addvspace{7pt plus 1pt}\normalfont}
  {\contentslabel{2.2em}}
  {\hspace*{-2.2em}}
  {\hfill\contentspage}
\titlecontents{section}[4.8em]
  {\normalfont\small}
  {\contentslabel{2.6em}}
  {\hspace*{-2.6em}}
  {\titlerule*[0.9em]{.}\contentspage}

%% ---------------------------------------------------------------------------
%% Index and bibliography: the index in two columns of small type; the bibliography
%% is "References", listed in the contents.
%% ---------------------------------------------------------------------------
% The last index page is balanced by multicol, which by default accepts a balanced column up
% to 12pt (a line) too long, so that one added entry can push a line into the bottom margin.
% With no overflow allowed, multicol fills that page normally and balances the rest on the
% next page instead; ragged columns keep balanced columns at their natural height.
\indexsetup{othercode={\small\raggedcolumns\maxbalancingoverflow=\z@}}
\@ifpackageloaded{natbib}{%
  \renewcommand{\bibsection}{\chapter*{References}%
    \addcontentsline{toc}{chapter}{References}}%
}{}

%% ---------------------------------------------------------------------------
%% Statements, examples and notes (amsthm). The three standard styles are redefined
%% with one generous spacing, so \theoremstyle{plain|definition|remark} in
%% main.tex picks up the design:
%%   plain       bold label, italic body    (Statement, Theorem, Lemma, ...)
%%   definition  bold label, upright body   (Definition, Example, Technical walkthrough)
%%   remark      italic label, upright body (Caution, Seminar question)
%% ---------------------------------------------------------------------------
\newlength{\trilogyenvskip}
\setlength{\trilogyenvskip}{10pt plus 3pt minus 2pt}
\newtheoremstyle{plain}{\trilogyenvskip}{\trilogyenvskip}
  {\itshape}{}{\bfseries}{.}{0.5em}{}
\newtheoremstyle{definition}{\trilogyenvskip}{\trilogyenvskip}
  {\normalfont}{}{\bfseries}{.}{0.5em}{}
\newtheoremstyle{remark}{\trilogyenvskip}{\trilogyenvskip}
  {\normalfont}{}{\itshape}{.}{0.5em}{}

%% Run-in notes ("Used in:", "Sources.", ...): an italic label and ordinary text,
%% a little space above. \trilogynote{<label>}{<text>}
\newcommand{\trilogynote}[2]{\par\addvspace{5pt plus 1pt}\noindent{\itshape #1}\enspace #2\par}

%% "Used in:" line after a statement: \usedin{<references>}
\newcommand{\usedin}[1]{\trilogynote{Used in:}{#1}}

%% Technical walkthrough: a numbered, multi-paragraph guided calculation,
%% \begin{walkthrough}[<optional title>] ... \end{walkthrough}. Definition style;
%% it ends with an open square at the right margin, like a proof, so the reader sees
%% where it stops. The name may be changed with \renewcommand{\walkthroughname}{...}.
\newcommand{\walkthroughname}{Technical walkthrough}
\newcommand{\walkthroughend}{\ensuremath{\square}}
%% Set-off statement: definition style, bold label and upright text, with the whole statement
%% indented 1.5em on both sides like the asides, instead of a box. Displays stay centred.
%% Declare with \theoremstyle{trilogysetoff}\newtheorem{<env>}{<name>}[chapter].
\newtheoremstyle{trilogysetoff}{\trilogyenvskip}{\trilogyenvskip}
  {\normalfont\leftskip=1.5em\rightskip=1.5em\relax}{}{\bfseries}{.}{0.5em}{}

\theoremstyle{definition}
\newtheorem{trilogywalk}{\walkthroughname}[chapter]
\newenvironment{walkthrough}[1][]
  {\ifx\relax#1\relax\begin{trilogywalk}\else\begin{trilogywalk}[#1]\fi
   \pushQED{\qed}\renewcommand{\qedsymbol}{\walkthroughend}}
  {\popQED\end{trilogywalk}}
\theoremstyle{plain}

%% Set-off asides, \begin{sayit} ... \end{sayit} ("How to say it") and
%% \begin{fiveminute} ... \end{fiveminute} ("Five-minute explanation"): a spaced
%% small-capitals label and a body indented on both sides; no box, no rule, no tint.
%% Inside sayit, \tophysicist and \tomathematician start the two voices.
\newenvironment{trilogyaside}[1]
  {\par\addvspace{11pt plus 3pt minus 2pt}%
   \list{}{\leftmargin=1.5em\rightmargin=1.5em\topsep=0pt\partopsep=0pt\parsep=0pt
     \listparindent=1.5em\itemsep=0pt}%
   \item\relax{\trilogy@spacedsc{#1}}\par\nopagebreak\vspace{3pt}\noindent\ignorespaces}
  {\endlist\par\addvspace{11pt plus 3pt minus 2pt}}
\newenvironment{sayit}{\begin{trilogyaside}{How to say it}}{\end{trilogyaside}}
\newenvironment{fiveminute}{\begin{trilogyaside}{Five-minute explanation}}{\end{trilogyaside}}
\newcommand{\tophysicist}{\par\addvspace{2pt}\noindent{\itshape To a physicist.}\enspace\ignorespaces}
\newcommand{\tomathematician}{\par\addvspace{4pt}\noindent{\itshape To a mathematician.}\enspace\ignorespaces}

%% Keeping a heading with what follows it: \trilogyneedspace{<length>} starts a new page at
%% once unless at least <length> of the current page remains. The test is made on the spot,
%% so it also works before a longtable (a breakpoint left for later would be taken under
%% longtable's output routine, which then repeats the table head above the heading).
\newcommand{\trilogyneedspace}[1]{\par\penalty-100\begingroup
  \setlength{\dimen@}{#1}\dimen@ii\pagegoal\advance\dimen@ii-\pagetotal
  \ifdim\dimen@>\dimen@ii\ifdim\dimen@ii>\z@\vfil\fi\break\fi\endgroup}

%% Captions: small type (11pt at the 12pt base), "Table A.1." with a full stop.
\usepackage[font=small,labelsep=period,skip=8pt]{caption}

%% Footnotes in \small rather than \footnotesize, for legibility on screen.
\usepackage{etoolbox}
\patchcmd{\@footnotetext}{\footnotesize}{\small}{}{}

%% ---------------------------------------------------------------------------
%% Wide reference tables. \widetablesetup, issued inside a group just before a
%% longtable, lets that table run \trilogy@overhang into each margin, centred on the
%% page: within the group \linewidth is the measure plus twice the overhang, so
%% p{<fraction>\linewidth} columns scale with it. With the 6.5in screen measure the
%% overhang is zero; the command is kept so that tables using it need no change.
%% ---------------------------------------------------------------------------
\newcommand{\trilogy@overhang}{0pt}
\newcommand{\widetablesetup}{%
  \setlength{\LTleft}{-\trilogy@overhang plus 1fill}%
  \setlength{\LTright}{-\trilogy@overhang plus 1fill}%
  \setlength{\linewidth}{\dimexpr\textwidth+2\dimexpr\trilogy@overhang\relax\relax}}

%% ---------------------------------------------------------------------------
%% Links: black and unboxed; they still work. No author in the PDF metadata.
%% The PDF opens on the document alone (no side panel), at full page width,
%% with the document title (not the file name) in the window bar.
%% ---------------------------------------------------------------------------
\usepackage[hidelinks,bookmarksnumbered=true]{hyperref}
\usepackage{bookmark}
\hypersetup{pdfauthor={},pdfpagemode=UseNone,pdfdisplaydoctitle=true,
  pdfstartview={FitH \hypercalcbp{\paperheight}}}

\makeatother
, after the class, fontenc and
%% inputenc, the AMS packages, natbib or biblatex, and imakeidx, and before the document's
%% theorem declarations and macros. The title, theorem names, local notes, and short
%% running heads are defined in main.tex.
%%
%% Principles: readability first, then understated elegance. The document is read mostly
%% on screen, so the page is set for that: 12pt type on a full-width letter page with 1in
%% side margins, one-sided, with the same quiet head on every page. Black text only; one
%% readable face with matching mathematics; headings distinguished by size, weight and
%% space rather than colour or rules; statements in the standard amsthm manner; classic
%% title page. The main.tex uses \documentclass[12pt,oneside,
%% openany]{book}.

\makeatletter

%% ---------------------------------------------------------------------------
%% Type: Palatino for text and mathematics (mathpazo, with true small capitals).
%% Its larger x-height and sturdier strokes read better than Latin Modern, on screen
%% and in dense formulas, and the math italic and symbols are drawn to match the text.
%% Latin Modern supplies the sans-serif and typewriter faces (URLs). bm comes after
%% the math fonts so that \bm finds bold Palatino.
%% ---------------------------------------------------------------------------
\usepackage{lmodern}
\usepackage[sc]{mathpazo}
\usepackage{bm}
\linespread{1.1}
\usepackage{microtype}
% Font expansion limited to 1.5% (default 2%): steadier colour, and it avoids the
% hairline overfull lines that maximal shrink plus protrusion can leave.
\microtypesetup{stretch=15,shrink=15}

%% ---------------------------------------------------------------------------
%% Page, set for reading on screen: US letter, 1in side margins (a 6.5in line, about
%% 85 characters of Palatino 12pt), and slightly smaller top and bottom margins with
%% the head and foot inside them, so a page fitted to the window shows large type.
%% ---------------------------------------------------------------------------
\usepackage{geometry}
\geometry{letterpaper,left=1in,right=1in,top=0.75in,bottom=0.75in,includeheadfoot,
  headheight=15pt,headsep=16pt,footskip=26pt}
\setlength{\parindent}{1.5em}
\setlength{\parskip}{0pt}
\setlength{\emergencystretch}{3em}
% A paragraph may set a slightly loose line rather than push a word into the margin.
\tolerance=1000
\clubpenalty=3000 \widowpenalty=3000 \displaywidowpenalty=1500
\raggedbottom
\allowdisplaybreaks[2]

%% Lists: compact, indented like a paragraph.
\usepackage{enumitem}
\setlist{topsep=4pt plus 2pt minus 1pt,itemsep=2pt plus 1pt,parsep=0pt}
\setlist[itemize]{leftmargin=1.5em}
\setlist[enumerate]{leftmargin=2em}

\usepackage[newparttoc]{titlesec}
\usepackage{titletoc,fancyhdr}

%% Letterspaced small capitals, for labels such as "Chapter 3".
\newcommand{\trilogy@spacedsc}[1]{{\scshape\textls[90]{\MakeLowercase{#1}}}}
%% Headings are ragged right and unhyphenated; a title that needs two lines is broken
%% into lines of similar length rather than leaving a single word on the second line.
\newcommand{\trilogy@headpar}{\rightskip=0pt plus 0.5\hsize\relax
  \parfillskip=0pt plus 0.5\hsize\relax\hyphenpenalty=10000\relax\exhyphenpenalty=10000\relax}

%% ---------------------------------------------------------------------------
%% Headings (titlesec): a regular-weight chapter title under a spaced small-capitals
%% label; bold section titles; plain part pages; no colour and no rules.
%% ---------------------------------------------------------------------------
\titleformat{\part}[display]
  {\normalfont\centering}
  {\large\trilogy@spacedsc{\partname\ \thepart}}
  {22pt}
  {\fontsize{22}{28}\selectfont}
\assignpagestyle{\part}{empty}

\titleformat{\chapter}[display]
  {\normalfont\trilogy@headpar}
  {\trilogy@spacedsc{\chaptertitlename\ \thechapter}}
  {12pt}
  {\fontsize{20}{25}\selectfont}
\titleformat{name=\chapter,numberless}[display]
  {\normalfont\trilogy@headpar}
  {}
  {0pt}
  {\fontsize{20}{25}\selectfont\trilogy@starmark}
\titlespacing*{\chapter}{0pt}{12pt}{30pt}
% An unnumbered chapter (front matter, contents, references, index) sets both running heads.
\newcommand{\trilogy@starmark}[1]{#1\markboth{#1}{#1}}

\titleformat{\section}
  {\normalfont\fontsize{13.5}{17}\selectfont\bfseries\trilogy@headpar}{\thesection}{0.75em}{}
\titlespacing*{\section}{0pt}{22pt plus 4pt minus 2pt}{9pt plus 2pt}
\titleformat{\subsection}
  {\normalfont\normalsize\bfseries\trilogy@headpar}{\thesubsection}{0.75em}{}
\titlespacing*{\subsection}{0pt}{15pt plus 3pt minus 2pt}{6pt plus 1pt}
\titleformat{\paragraph}[runin]
  {\normalfont\normalsize\bfseries}{}{0pt}{}
\titlespacing*{\paragraph}{0pt}{9pt plus 2pt minus 1pt}{0.6em}

\setcounter{tocdepth}{1}
\setcounter{secnumdepth}{1}

%% ---------------------------------------------------------------------------
%% Running heads (one-sided, the same on every page): the chapter number and title at
%% the left, and at the right the section number and title followed by the page number,
%% all in small italic except the numbers, with no rule. When chapter and section titles
%% do not both fit on the line, the right head shows only the section number (§1.2). Chapter-
%% opening pages carry only a centred page number at the foot; part pages carry nothing.
%%
%% A long title can be given a short form in main.tex:
%%   \shorthead{<title exactly as in the source>}{<short form>}
%% A chapter head too wide for the line falls back to "Chapter N" (with a warning).
%% ---------------------------------------------------------------------------
\newcommand{\shorthead}[2]{\expandafter\def\csname trilogy@sh@\detokenize{#1}\endcsname{#2}}
\newcommand{\trilogy@lookup}[1]{%
  \ifcsname trilogy@sh@\detokenize{#1}\endcsname
    \expandafter\let\expandafter\trilogy@headtext\csname trilogy@sh@\detokenize{#1}\endcsname
  \else
    \def\trilogy@headtext{#1}%
  \fi}
\newcommand{\trilogy@headfont}{\normalfont\small\itshape}
% A mark is \trilogy@hd{number}{title}; the number is empty for unnumbered chapters.
\DeclareRobustCommand{\trilogy@hd}[2]{%
  \ifx\relax#1\relax\else{\upshape#1}\hspace{0.8em}\fi#2}
\newsavebox{\trilogy@chapbox}
\newsavebox{\trilogy@secbox}
\newsavebox{\trilogy@cachedchapbox}
\newsavebox{\trilogy@cachedsecbox}
\newif\iftrilogy@headtoolong
\newif\iftrilogy@cachedheadtoolong
\let\trilogy@cachedheadkey\relax
% Compute the two boxes only when their marks or dimensions change.
\newcommand{\trilogy@computeheads}{%
  \trilogy@headtoolongfalse
  \sbox\trilogy@chapbox{\trilogy@headfont\nouppercase{\leftmark}}%
  \sbox\trilogy@secbox{}%
  \protected@edef\trilogy@l{\leftmark}%
  \protected@edef\trilogy@r{\rightmark}%
  \ifx\trilogy@r\@empty\else\ifx\trilogy@r\trilogy@l\else
    \sbox\trilogy@secbox{\trilogy@headfont\nouppercase{\rightmark}}%
  \fi\fi
  \dimen@=\wd\trilogy@chapbox\advance\dimen@ 4em\relax
  \ifdim\dimen@>\headwidth
    \sbox\trilogy@chapbox{\trilogy@headfont\@chapapp\ \thechapter}%
    \trilogy@headtoolongtrue
  \fi
  \dimen@=\wd\trilogy@chapbox\advance\dimen@\wd\trilogy@secbox\advance\dimen@ 8em\relax
  \ifdim\dimen@>\headwidth
    \sbox\trilogy@secbox{\trilogy@headfont
      \expandafter\let\csname trilogy@hd \endcsname\trilogy@hdnum\nouppercase{\rightmark}}%
  \fi}
% Both fancyhdr fields, and successive pages with the same marks, reuse these
% boxes. The cache is global because fancyhdr evaluates its fields in groups;
% the page number and page-specific warnings remain outside the cached boxes.
\newcommand{\trilogy@heads}[1]{%
  \protected@edef\trilogy@headkey{%
    {\leftmark}{\rightmark}{\the\headwidth}%
    {\fontname\font}{\the\dimexpr1em\relax}{\thechapter}{\@chapapp}}%
  \ifx\trilogy@headkey\trilogy@cachedheadkey\else
    \trilogy@computeheads
    \global\setbox\trilogy@cachedchapbox=\copy\trilogy@chapbox
    \global\setbox\trilogy@cachedsecbox=\copy\trilogy@secbox
    \global\let\trilogy@cachedheadkey\trilogy@headkey
    \iftrilogy@headtoolong
      \global\trilogy@cachedheadtoolongtrue
    \else
      \global\trilogy@cachedheadtoolongfalse
    \fi
  \fi
  \setbox\trilogy@chapbox=\copy\trilogy@cachedchapbox
  \setbox\trilogy@secbox=\copy\trilogy@cachedsecbox
  \iftrilogy@cachedheadtoolong
    \ifx\relax#1\relax\else\PackageWarningNoLine{trilogy-style}{Running head too long on
      page \thepage; give the chapter a \string\shorthead}\fi
  \fi}
% The section number alone, for a right head with no room for the section title. (A mark
% holds the robust command's inner macro "\trilogy@hd ", which is what is replaced above.)
\newcommand{\trilogy@hdnum}[2]{\ifx\relax#1\relax\else\S\,{\upshape#1}\fi}

\pagestyle{fancy}
\fancyhf{}
\fancyhead[L]{\trilogy@heads{warn}\usebox\trilogy@chapbox}
\fancyhead[R]{\trilogy@heads{}%
  \ifdim\wd\trilogy@secbox>\z@\usebox\trilogy@secbox\hspace{2em}\fi
  \normalfont\small\thepage}
\renewcommand{\headrulewidth}{0pt}
\renewcommand{\footrulewidth}{0pt}
\fancypagestyle{plain}{\fancyhf{}\fancyfoot[C]{\normalfont\small\thepage}%
  \renewcommand{\headrulewidth}{0pt}\renewcommand{\footrulewidth}{0pt}}
\renewcommand{\chaptermark}[1]{\trilogy@lookup{#1}%
  \markboth{\trilogy@hd{\ifnum\c@secnumdepth>\m@ne\if@mainmatter\thechapter\fi\fi}{\trilogy@headtext}}{}}
\renewcommand{\sectionmark}[1]{\trilogy@lookup{#1}%
  \markright{\trilogy@hd{\ifnum\c@secnumdepth>\z@\thesection\fi}{\trilogy@headtext}}}

%% ---------------------------------------------------------------------------
%% Table of contents (titletoc): parts in spaced small capitals, chapters in roman with
%% hanging numbers, sections indented and smaller with sparse leaders.
%% ---------------------------------------------------------------------------
\contentsmargin{2.2em}
\titlecontents{part}[0pt]
  {\addvspace{16pt plus 2pt}\normalfont}
  {\trilogy@spacedsc{\partname\ \thecontentslabel}\hspace{1em}}
  {}
  {\hfill\contentspage}
  [\addvspace{4pt}]
\titlecontents{chapter}[2.2em]
  {\addvspace{7pt plus 1pt}\normalfont}
  {\contentslabel{2.2em}}
  {\hspace*{-2.2em}}
  {\hfill\contentspage}
\titlecontents{section}[4.8em]
  {\normalfont\small}
  {\contentslabel{2.6em}}
  {\hspace*{-2.6em}}
  {\titlerule*[0.9em]{.}\contentspage}

%% ---------------------------------------------------------------------------
%% Index and bibliography: the index in two columns of small type; the bibliography
%% is "References", listed in the contents.
%% ---------------------------------------------------------------------------
% The last index page is balanced by multicol, which by default accepts a balanced column up
% to 12pt (a line) too long, so that one added entry can push a line into the bottom margin.
% With no overflow allowed, multicol fills that page normally and balances the rest on the
% next page instead; ragged columns keep balanced columns at their natural height.
\indexsetup{othercode={\small\raggedcolumns\maxbalancingoverflow=\z@}}
\@ifpackageloaded{natbib}{%
  \renewcommand{\bibsection}{\chapter*{References}%
    \addcontentsline{toc}{chapter}{References}}%
}{}

%% ---------------------------------------------------------------------------
%% Statements, examples and notes (amsthm). The three standard styles are redefined
%% with one generous spacing, so \theoremstyle{plain|definition|remark} in
%% main.tex picks up the design:
%%   plain       bold label, italic body    (Statement, Theorem, Lemma, ...)
%%   definition  bold label, upright body   (Definition, Example, Technical walkthrough)
%%   remark      italic label, upright body (Caution, Seminar question)
%% ---------------------------------------------------------------------------
\newlength{\trilogyenvskip}
\setlength{\trilogyenvskip}{10pt plus 3pt minus 2pt}
\newtheoremstyle{plain}{\trilogyenvskip}{\trilogyenvskip}
  {\itshape}{}{\bfseries}{.}{0.5em}{}
\newtheoremstyle{definition}{\trilogyenvskip}{\trilogyenvskip}
  {\normalfont}{}{\bfseries}{.}{0.5em}{}
\newtheoremstyle{remark}{\trilogyenvskip}{\trilogyenvskip}
  {\normalfont}{}{\itshape}{.}{0.5em}{}

%% Run-in notes ("Used in:", "Sources.", ...): an italic label and ordinary text,
%% a little space above. \trilogynote{<label>}{<text>}
\newcommand{\trilogynote}[2]{\par\addvspace{5pt plus 1pt}\noindent{\itshape #1}\enspace #2\par}

%% "Used in:" line after a statement: \usedin{<references>}
\newcommand{\usedin}[1]{\trilogynote{Used in:}{#1}}

%% Technical walkthrough: a numbered, multi-paragraph guided calculation,
%% \begin{walkthrough}[<optional title>] ... \end{walkthrough}. Definition style;
%% it ends with an open square at the right margin, like a proof, so the reader sees
%% where it stops. The name may be changed with \renewcommand{\walkthroughname}{...}.
\newcommand{\walkthroughname}{Technical walkthrough}
\newcommand{\walkthroughend}{\ensuremath{\square}}
%% Set-off statement: definition style, bold label and upright text, with the whole statement
%% indented 1.5em on both sides like the asides, instead of a box. Displays stay centred.
%% Declare with \theoremstyle{trilogysetoff}\newtheorem{<env>}{<name>}[chapter].
\newtheoremstyle{trilogysetoff}{\trilogyenvskip}{\trilogyenvskip}
  {\normalfont\leftskip=1.5em\rightskip=1.5em\relax}{}{\bfseries}{.}{0.5em}{}

\theoremstyle{definition}
\newtheorem{trilogywalk}{\walkthroughname}[chapter]
\newenvironment{walkthrough}[1][]
  {\ifx\relax#1\relax\begin{trilogywalk}\else\begin{trilogywalk}[#1]\fi
   \pushQED{\qed}\renewcommand{\qedsymbol}{\walkthroughend}}
  {\popQED\end{trilogywalk}}
\theoremstyle{plain}

%% Set-off asides, \begin{sayit} ... \end{sayit} ("How to say it") and
%% \begin{fiveminute} ... \end{fiveminute} ("Five-minute explanation"): a spaced
%% small-capitals label and a body indented on both sides; no box, no rule, no tint.
%% Inside sayit, \tophysicist and \tomathematician start the two voices.
\newenvironment{trilogyaside}[1]
  {\par\addvspace{11pt plus 3pt minus 2pt}%
   \list{}{\leftmargin=1.5em\rightmargin=1.5em\topsep=0pt\partopsep=0pt\parsep=0pt
     \listparindent=1.5em\itemsep=0pt}%
   \item\relax{\trilogy@spacedsc{#1}}\par\nopagebreak\vspace{3pt}\noindent\ignorespaces}
  {\endlist\par\addvspace{11pt plus 3pt minus 2pt}}
\newenvironment{sayit}{\begin{trilogyaside}{How to say it}}{\end{trilogyaside}}
\newenvironment{fiveminute}{\begin{trilogyaside}{Five-minute explanation}}{\end{trilogyaside}}
\newcommand{\tophysicist}{\par\addvspace{2pt}\noindent{\itshape To a physicist.}\enspace\ignorespaces}
\newcommand{\tomathematician}{\par\addvspace{4pt}\noindent{\itshape To a mathematician.}\enspace\ignorespaces}

%% Keeping a heading with what follows it: \trilogyneedspace{<length>} starts a new page at
%% once unless at least <length> of the current page remains. The test is made on the spot,
%% so it also works before a longtable (a breakpoint left for later would be taken under
%% longtable's output routine, which then repeats the table head above the heading).
\newcommand{\trilogyneedspace}[1]{\par\penalty-100\begingroup
  \setlength{\dimen@}{#1}\dimen@ii\pagegoal\advance\dimen@ii-\pagetotal
  \ifdim\dimen@>\dimen@ii\ifdim\dimen@ii>\z@\vfil\fi\break\fi\endgroup}

%% Captions: small type (11pt at the 12pt base), "Table A.1." with a full stop.
\usepackage[font=small,labelsep=period,skip=8pt]{caption}

%% Footnotes in \small rather than \footnotesize, for legibility on screen.
\usepackage{etoolbox}
\patchcmd{\@footnotetext}{\footnotesize}{\small}{}{}

%% ---------------------------------------------------------------------------
%% Wide reference tables. \widetablesetup, issued inside a group just before a
%% longtable, lets that table run \trilogy@overhang into each margin, centred on the
%% page: within the group \linewidth is the measure plus twice the overhang, so
%% p{<fraction>\linewidth} columns scale with it. With the 6.5in screen measure the
%% overhang is zero; the command is kept so that tables using it need no change.
%% ---------------------------------------------------------------------------
\newcommand{\trilogy@overhang}{0pt}
\newcommand{\widetablesetup}{%
  \setlength{\LTleft}{-\trilogy@overhang plus 1fill}%
  \setlength{\LTright}{-\trilogy@overhang plus 1fill}%
  \setlength{\linewidth}{\dimexpr\textwidth+2\dimexpr\trilogy@overhang\relax\relax}}

%% ---------------------------------------------------------------------------
%% Links: black and unboxed; they still work. No author in the PDF metadata.
%% The PDF opens on the document alone (no side panel), at full page width,
%% with the document title (not the file name) in the window bar.
%% ---------------------------------------------------------------------------
\usepackage[hidelinks,bookmarksnumbered=true]{hyperref}
\usepackage{bookmark}
\hypersetup{pdfauthor={},pdfpagemode=UseNone,pdfdisplaydoctitle=true,
  pdfstartview={FitH \hypercalcbp{\paperheight}}}

\makeatother
, after the class, fontenc and
%% inputenc, the AMS packages, natbib or biblatex, and imakeidx, and before the document's
%% theorem declarations and macros. The title, theorem names, local notes, and short
%% running heads are defined in main.tex.
%%
%% Principles: readability first, then understated elegance. The document is read mostly
%% on screen, so the page is set for that: 12pt type on a full-width letter page with 1in
%% side margins, one-sided, with the same quiet head on every page. Black text only; one
%% readable face with matching mathematics; headings distinguished by size, weight and
%% space rather than colour or rules; statements in the standard amsthm manner; classic
%% title page. The main.tex uses \documentclass[12pt,oneside,
%% openany]{book}.

\makeatletter

%% ---------------------------------------------------------------------------
%% Type: Palatino for text and mathematics (mathpazo, with true small capitals).
%% Its larger x-height and sturdier strokes read better than Latin Modern, on screen
%% and in dense formulas, and the math italic and symbols are drawn to match the text.
%% Latin Modern supplies the sans-serif and typewriter faces (URLs). bm comes after
%% the math fonts so that \bm finds bold Palatino.
%% ---------------------------------------------------------------------------
\usepackage{lmodern}
\usepackage[sc]{mathpazo}
\usepackage{bm}
\linespread{1.1}
\usepackage{microtype}
% Font expansion limited to 1.5% (default 2%): steadier colour, and it avoids the
% hairline overfull lines that maximal shrink plus protrusion can leave.
\microtypesetup{stretch=15,shrink=15}

%% ---------------------------------------------------------------------------
%% Page, set for reading on screen: US letter, 1in side margins (a 6.5in line, about
%% 85 characters of Palatino 12pt), and slightly smaller top and bottom margins with
%% the head and foot inside them, so a page fitted to the window shows large type.
%% ---------------------------------------------------------------------------
\usepackage{geometry}
\geometry{letterpaper,left=1in,right=1in,top=0.75in,bottom=0.75in,includeheadfoot,
  headheight=15pt,headsep=16pt,footskip=26pt}
\setlength{\parindent}{1.5em}
\setlength{\parskip}{0pt}
\setlength{\emergencystretch}{3em}
% A paragraph may set a slightly loose line rather than push a word into the margin.
\tolerance=1000
\clubpenalty=3000 \widowpenalty=3000 \displaywidowpenalty=1500
\raggedbottom
\allowdisplaybreaks[2]

%% Lists: compact, indented like a paragraph.
\usepackage{enumitem}
\setlist{topsep=4pt plus 2pt minus 1pt,itemsep=2pt plus 1pt,parsep=0pt}
\setlist[itemize]{leftmargin=1.5em}
\setlist[enumerate]{leftmargin=2em}

\usepackage[newparttoc]{titlesec}
\usepackage{titletoc,fancyhdr}

%% Letterspaced small capitals, for labels such as "Chapter 3".
\newcommand{\trilogy@spacedsc}[1]{{\scshape\textls[90]{\MakeLowercase{#1}}}}
%% Headings are ragged right and unhyphenated; a title that needs two lines is broken
%% into lines of similar length rather than leaving a single word on the second line.
\newcommand{\trilogy@headpar}{\rightskip=0pt plus 0.5\hsize\relax
  \parfillskip=0pt plus 0.5\hsize\relax\hyphenpenalty=10000\relax\exhyphenpenalty=10000\relax}

%% ---------------------------------------------------------------------------
%% Headings (titlesec): a regular-weight chapter title under a spaced small-capitals
%% label; bold section titles; plain part pages; no colour and no rules.
%% ---------------------------------------------------------------------------
\titleformat{\part}[display]
  {\normalfont\centering}
  {\large\trilogy@spacedsc{\partname\ \thepart}}
  {22pt}
  {\fontsize{22}{28}\selectfont}
\assignpagestyle{\part}{empty}

\titleformat{\chapter}[display]
  {\normalfont\trilogy@headpar}
  {\trilogy@spacedsc{\chaptertitlename\ \thechapter}}
  {12pt}
  {\fontsize{20}{25}\selectfont}
\titleformat{name=\chapter,numberless}[display]
  {\normalfont\trilogy@headpar}
  {}
  {0pt}
  {\fontsize{20}{25}\selectfont\trilogy@starmark}
\titlespacing*{\chapter}{0pt}{12pt}{30pt}
% An unnumbered chapter (front matter, contents, references, index) sets both running heads.
\newcommand{\trilogy@starmark}[1]{#1\markboth{#1}{#1}}

\titleformat{\section}
  {\normalfont\fontsize{13.5}{17}\selectfont\bfseries\trilogy@headpar}{\thesection}{0.75em}{}
\titlespacing*{\section}{0pt}{22pt plus 4pt minus 2pt}{9pt plus 2pt}
\titleformat{\subsection}
  {\normalfont\normalsize\bfseries\trilogy@headpar}{\thesubsection}{0.75em}{}
\titlespacing*{\subsection}{0pt}{15pt plus 3pt minus 2pt}{6pt plus 1pt}
\titleformat{\paragraph}[runin]
  {\normalfont\normalsize\bfseries}{}{0pt}{}
\titlespacing*{\paragraph}{0pt}{9pt plus 2pt minus 1pt}{0.6em}

\setcounter{tocdepth}{1}
\setcounter{secnumdepth}{1}

%% ---------------------------------------------------------------------------
%% Running heads (one-sided, the same on every page): the chapter number and title at
%% the left, and at the right the section number and title followed by the page number,
%% all in small italic except the numbers, with no rule. When chapter and section titles
%% do not both fit on the line, the right head shows only the section number (§1.2). Chapter-
%% opening pages carry only a centred page number at the foot; part pages carry nothing.
%%
%% A long title can be given a short form in main.tex:
%%   \shorthead{<title exactly as in the source>}{<short form>}
%% A chapter head too wide for the line falls back to "Chapter N" (with a warning).
%% ---------------------------------------------------------------------------
\newcommand{\shorthead}[2]{\expandafter\def\csname trilogy@sh@\detokenize{#1}\endcsname{#2}}
\newcommand{\trilogy@lookup}[1]{%
  \ifcsname trilogy@sh@\detokenize{#1}\endcsname
    \expandafter\let\expandafter\trilogy@headtext\csname trilogy@sh@\detokenize{#1}\endcsname
  \else
    \def\trilogy@headtext{#1}%
  \fi}
\newcommand{\trilogy@headfont}{\normalfont\small\itshape}
% A mark is \trilogy@hd{number}{title}; the number is empty for unnumbered chapters.
\DeclareRobustCommand{\trilogy@hd}[2]{%
  \ifx\relax#1\relax\else{\upshape#1}\hspace{0.8em}\fi#2}
\newsavebox{\trilogy@chapbox}
\newsavebox{\trilogy@secbox}
\newsavebox{\trilogy@cachedchapbox}
\newsavebox{\trilogy@cachedsecbox}
\newif\iftrilogy@headtoolong
\newif\iftrilogy@cachedheadtoolong
\let\trilogy@cachedheadkey\relax
% Compute the two boxes only when their marks or dimensions change.
\newcommand{\trilogy@computeheads}{%
  \trilogy@headtoolongfalse
  \sbox\trilogy@chapbox{\trilogy@headfont\nouppercase{\leftmark}}%
  \sbox\trilogy@secbox{}%
  \protected@edef\trilogy@l{\leftmark}%
  \protected@edef\trilogy@r{\rightmark}%
  \ifx\trilogy@r\@empty\else\ifx\trilogy@r\trilogy@l\else
    \sbox\trilogy@secbox{\trilogy@headfont\nouppercase{\rightmark}}%
  \fi\fi
  \dimen@=\wd\trilogy@chapbox\advance\dimen@ 4em\relax
  \ifdim\dimen@>\headwidth
    \sbox\trilogy@chapbox{\trilogy@headfont\@chapapp\ \thechapter}%
    \trilogy@headtoolongtrue
  \fi
  \dimen@=\wd\trilogy@chapbox\advance\dimen@\wd\trilogy@secbox\advance\dimen@ 8em\relax
  \ifdim\dimen@>\headwidth
    \sbox\trilogy@secbox{\trilogy@headfont
      \expandafter\let\csname trilogy@hd \endcsname\trilogy@hdnum\nouppercase{\rightmark}}%
  \fi}
% Both fancyhdr fields, and successive pages with the same marks, reuse these
% boxes. The cache is global because fancyhdr evaluates its fields in groups;
% the page number and page-specific warnings remain outside the cached boxes.
\newcommand{\trilogy@heads}[1]{%
  \protected@edef\trilogy@headkey{%
    {\leftmark}{\rightmark}{\the\headwidth}%
    {\fontname\font}{\the\dimexpr1em\relax}{\thechapter}{\@chapapp}}%
  \ifx\trilogy@headkey\trilogy@cachedheadkey\else
    \trilogy@computeheads
    \global\setbox\trilogy@cachedchapbox=\copy\trilogy@chapbox
    \global\setbox\trilogy@cachedsecbox=\copy\trilogy@secbox
    \global\let\trilogy@cachedheadkey\trilogy@headkey
    \iftrilogy@headtoolong
      \global\trilogy@cachedheadtoolongtrue
    \else
      \global\trilogy@cachedheadtoolongfalse
    \fi
  \fi
  \setbox\trilogy@chapbox=\copy\trilogy@cachedchapbox
  \setbox\trilogy@secbox=\copy\trilogy@cachedsecbox
  \iftrilogy@cachedheadtoolong
    \ifx\relax#1\relax\else\PackageWarningNoLine{trilogy-style}{Running head too long on
      page \thepage; give the chapter a \string\shorthead}\fi
  \fi}
% The section number alone, for a right head with no room for the section title. (A mark
% holds the robust command's inner macro "\trilogy@hd ", which is what is replaced above.)
\newcommand{\trilogy@hdnum}[2]{\ifx\relax#1\relax\else\S\,{\upshape#1}\fi}

\pagestyle{fancy}
\fancyhf{}
\fancyhead[L]{\trilogy@heads{warn}\usebox\trilogy@chapbox}
\fancyhead[R]{\trilogy@heads{}%
  \ifdim\wd\trilogy@secbox>\z@\usebox\trilogy@secbox\hspace{2em}\fi
  \normalfont\small\thepage}
\renewcommand{\headrulewidth}{0pt}
\renewcommand{\footrulewidth}{0pt}
\fancypagestyle{plain}{\fancyhf{}\fancyfoot[C]{\normalfont\small\thepage}%
  \renewcommand{\headrulewidth}{0pt}\renewcommand{\footrulewidth}{0pt}}
\renewcommand{\chaptermark}[1]{\trilogy@lookup{#1}%
  \markboth{\trilogy@hd{\ifnum\c@secnumdepth>\m@ne\if@mainmatter\thechapter\fi\fi}{\trilogy@headtext}}{}}
\renewcommand{\sectionmark}[1]{\trilogy@lookup{#1}%
  \markright{\trilogy@hd{\ifnum\c@secnumdepth>\z@\thesection\fi}{\trilogy@headtext}}}

%% ---------------------------------------------------------------------------
%% Table of contents (titletoc): parts in spaced small capitals, chapters in roman with
%% hanging numbers, sections indented and smaller with sparse leaders.
%% ---------------------------------------------------------------------------
\contentsmargin{2.2em}
\titlecontents{part}[0pt]
  {\addvspace{16pt plus 2pt}\normalfont}
  {\trilogy@spacedsc{\partname\ \thecontentslabel}\hspace{1em}}
  {}
  {\hfill\contentspage}
  [\addvspace{4pt}]
\titlecontents{chapter}[2.2em]
  {\addvspace{7pt plus 1pt}\normalfont}
  {\contentslabel{2.2em}}
  {\hspace*{-2.2em}}
  {\hfill\contentspage}
\titlecontents{section}[4.8em]
  {\normalfont\small}
  {\contentslabel{2.6em}}
  {\hspace*{-2.6em}}
  {\titlerule*[0.9em]{.}\contentspage}

%% ---------------------------------------------------------------------------
%% Index and bibliography: the index in two columns of small type; the bibliography
%% is "References", listed in the contents.
%% ---------------------------------------------------------------------------
% The last index page is balanced by multicol, which by default accepts a balanced column up
% to 12pt (a line) too long, so that one added entry can push a line into the bottom margin.
% With no overflow allowed, multicol fills that page normally and balances the rest on the
% next page instead; ragged columns keep balanced columns at their natural height.
\indexsetup{othercode={\small\raggedcolumns\maxbalancingoverflow=\z@}}
\@ifpackageloaded{natbib}{%
  \renewcommand{\bibsection}{\chapter*{References}%
    \addcontentsline{toc}{chapter}{References}}%
}{}

%% ---------------------------------------------------------------------------
%% Statements, examples and notes (amsthm). The three standard styles are redefined
%% with one generous spacing, so \theoremstyle{plain|definition|remark} in
%% main.tex picks up the design:
%%   plain       bold label, italic body    (Statement, Theorem, Lemma, ...)
%%   definition  bold label, upright body   (Definition, Example, Technical walkthrough)
%%   remark      italic label, upright body (Caution, Seminar question)
%% ---------------------------------------------------------------------------
\newlength{\trilogyenvskip}
\setlength{\trilogyenvskip}{10pt plus 3pt minus 2pt}
\newtheoremstyle{plain}{\trilogyenvskip}{\trilogyenvskip}
  {\itshape}{}{\bfseries}{.}{0.5em}{}
\newtheoremstyle{definition}{\trilogyenvskip}{\trilogyenvskip}
  {\normalfont}{}{\bfseries}{.}{0.5em}{}
\newtheoremstyle{remark}{\trilogyenvskip}{\trilogyenvskip}
  {\normalfont}{}{\itshape}{.}{0.5em}{}

%% Run-in notes ("Used in:", "Sources.", ...): an italic label and ordinary text,
%% a little space above. \trilogynote{<label>}{<text>}
\newcommand{\trilogynote}[2]{\par\addvspace{5pt plus 1pt}\noindent{\itshape #1}\enspace #2\par}

%% "Used in:" line after a statement: \usedin{<references>}
\newcommand{\usedin}[1]{\trilogynote{Used in:}{#1}}

%% Technical walkthrough: a numbered, multi-paragraph guided calculation,
%% \begin{walkthrough}[<optional title>] ... \end{walkthrough}. Definition style;
%% it ends with an open square at the right margin, like a proof, so the reader sees
%% where it stops. The name may be changed with \renewcommand{\walkthroughname}{...}.
\newcommand{\walkthroughname}{Technical walkthrough}
\newcommand{\walkthroughend}{\ensuremath{\square}}
%% Set-off statement: definition style, bold label and upright text, with the whole statement
%% indented 1.5em on both sides like the asides, instead of a box. Displays stay centred.
%% Declare with \theoremstyle{trilogysetoff}\newtheorem{<env>}{<name>}[chapter].
\newtheoremstyle{trilogysetoff}{\trilogyenvskip}{\trilogyenvskip}
  {\normalfont\leftskip=1.5em\rightskip=1.5em\relax}{}{\bfseries}{.}{0.5em}{}

\theoremstyle{definition}
\newtheorem{trilogywalk}{\walkthroughname}[chapter]
\newenvironment{walkthrough}[1][]
  {\ifx\relax#1\relax\begin{trilogywalk}\else\begin{trilogywalk}[#1]\fi
   \pushQED{\qed}\renewcommand{\qedsymbol}{\walkthroughend}}
  {\popQED\end{trilogywalk}}
\theoremstyle{plain}

%% Set-off asides, \begin{sayit} ... \end{sayit} ("How to say it") and
%% \begin{fiveminute} ... \end{fiveminute} ("Five-minute explanation"): a spaced
%% small-capitals label and a body indented on both sides; no box, no rule, no tint.
%% Inside sayit, \tophysicist and \tomathematician start the two voices.
\newenvironment{trilogyaside}[1]
  {\par\addvspace{11pt plus 3pt minus 2pt}%
   \list{}{\leftmargin=1.5em\rightmargin=1.5em\topsep=0pt\partopsep=0pt\parsep=0pt
     \listparindent=1.5em\itemsep=0pt}%
   \item\relax{\trilogy@spacedsc{#1}}\par\nopagebreak\vspace{3pt}\noindent\ignorespaces}
  {\endlist\par\addvspace{11pt plus 3pt minus 2pt}}
\newenvironment{sayit}{\begin{trilogyaside}{How to say it}}{\end{trilogyaside}}
\newenvironment{fiveminute}{\begin{trilogyaside}{Five-minute explanation}}{\end{trilogyaside}}
\newcommand{\tophysicist}{\par\addvspace{2pt}\noindent{\itshape To a physicist.}\enspace\ignorespaces}
\newcommand{\tomathematician}{\par\addvspace{4pt}\noindent{\itshape To a mathematician.}\enspace\ignorespaces}

%% Keeping a heading with what follows it: \trilogyneedspace{<length>} starts a new page at
%% once unless at least <length> of the current page remains. The test is made on the spot,
%% so it also works before a longtable (a breakpoint left for later would be taken under
%% longtable's output routine, which then repeats the table head above the heading).
\newcommand{\trilogyneedspace}[1]{\par\penalty-100\begingroup
  \setlength{\dimen@}{#1}\dimen@ii\pagegoal\advance\dimen@ii-\pagetotal
  \ifdim\dimen@>\dimen@ii\ifdim\dimen@ii>\z@\vfil\fi\break\fi\endgroup}

%% Captions: small type (11pt at the 12pt base), "Table A.1." with a full stop.
\usepackage[font=small,labelsep=period,skip=8pt]{caption}

%% Footnotes in \small rather than \footnotesize, for legibility on screen.
\usepackage{etoolbox}
\patchcmd{\@footnotetext}{\footnotesize}{\small}{}{}

%% ---------------------------------------------------------------------------
%% Wide reference tables. \widetablesetup, issued inside a group just before a
%% longtable, lets that table run \trilogy@overhang into each margin, centred on the
%% page: within the group \linewidth is the measure plus twice the overhang, so
%% p{<fraction>\linewidth} columns scale with it. With the 6.5in screen measure the
%% overhang is zero; the command is kept so that tables using it need no change.
%% ---------------------------------------------------------------------------
\newcommand{\trilogy@overhang}{0pt}
\newcommand{\widetablesetup}{%
  \setlength{\LTleft}{-\trilogy@overhang plus 1fill}%
  \setlength{\LTright}{-\trilogy@overhang plus 1fill}%
  \setlength{\linewidth}{\dimexpr\textwidth+2\dimexpr\trilogy@overhang\relax\relax}}

%% ---------------------------------------------------------------------------
%% Links: black and unboxed; they still work. No author in the PDF metadata.
%% The PDF opens on the document alone (no side panel), at full page width,
%% with the document title (not the file name) in the window bar.
%% ---------------------------------------------------------------------------
\usepackage[hidelinks,bookmarksnumbered=true]{hyperref}
\usepackage{bookmark}
\hypersetup{pdfauthor={},pdfpagemode=UseNone,pdfdisplaydoctitle=true,
  pdfstartview={FitH \hypercalcbp{\paperheight}}}

\makeatother
